\section{Related Work}
\label{sec:related}

\paragraph{Fault-diagnosis benchmarks and standardization.}
The closest domain work is the application-oriented benchmark of Zhao, Zio,
and Shen~\cite{zhao2024dgbenchmark}, which unifies public rotating-machinery
datasets under a common domain-generalization protocol with open code
(\texttt{DG\_benchmark}~\cite{dgbenchmarkrepo}), and a 2025 standardization
framework in IEEE Access that formalizes fault-level and operating-condition
factors for controlled evaluation~\cite{ieeeaccess2025standard}. Both respond
to documented comparability failures of ad hoc setups, and both standardize
the \emph{content} of evaluation: datasets, shift factors, partitions,
metrics. Neither prescribes, enforces, or verifies the \emph{validity
mechanisms} of the experiments built on top---recording-level leakage
control, test isolation, seed completeness, or provenance chains. That layer
is the subject of this paper, and the two lines are complementary: the
protocols those works define could be expressed as EdgeFault-Bench task
manifests and would then inherit the kernel's enforcement.

\paragraph{Domain-generalization evaluation practice.}
DomainBed~\cite{gulrajani2021domainbed} established rigorous
model-selection discipline for domain generalization in vision, showing that
fair selection changes algorithm rankings. Its concern---that evaluation
protocol, not algorithm, often drives conclusions---is the same concern we
transport into industrial experimentation; DomainBed, however, distributes
datasets and training scripts rather than installable integrity contracts,
and its machine-readable artefacts do not model machinery-specific notions
such as the recording as the unit of independence. General time-series
ecosystems such as aeon~\cite{middlehurst2025aeon} provide algorithm
breadth; EdgeFault-Bench deliberately reuses such ecosystems for estimators
where possible and concentrates on what they do not provide:
condition-held-out task semantics, source-verified streaming, and audit
artefacts.

\paragraph{Experiment infrastructure and reproducibility.}
MLflow~\cite{mlflow} and similar trackers log whatever an experiment emits;
they neither define partitions nor reject a bundle with a missing seed.
Recent reproducibility analyses argue that releasing code and dependency
lists alone does not make ML benchmarks runnable or
trustworthy~\cite{lyu2026reproducibility}, and infrastructure benchmarking
has begun treating the experimental platform itself as the object of
study~\cite{kramer2025mlops}; benchmark suites for autonomous research
systems evaluate whether agents can replicate results~\cite{mlreplicate2026}.
These works share our diagnosis that the platform layer decides
reproducibility. Our contribution to that conversation is a platform whose
integrity properties are \emph{specified formally and attacked deliberately}
rather than asserted, with the fault-injection results to show for it.

\paragraph{Formal methods in ML tooling.}
TLA\textsuperscript{+} and TLC are established tools for specifying and
model-checking concurrent and distributed systems~\cite{lamport2002specifying};
property-based testing with Hypothesis~\cite{hypothesis} is an established
implementation-level verification technique. To our knowledge, applying
this stack---a guarded lifecycle state machine, model-checked invariants,
oracle-based property tests, and a mutation catalogue---to the integrity of
ML \emph{benchmark infrastructure} has not been done before; adjacent
formal-methods work targets training systems and distributed learning, not
the experiment lifecycle itself. A dedicated literature pass before
submission will re-examine this positioning.
